\exercisesection{2}

\begin{enumerate}
\def\labelenumi{\arabic{enumi}.}
\tightlist
\item
  乘性子集 \(S\) （环 \(A\) 的）称为饱和的，如果关系 \(xy \in S\) 蕴含 \(x \in S\) 且 \(y \in S\) 。
\end{enumerate}

\begin{enumerate}
\def\labelenumi{(\alph{enumi})}
\item
  为使子集 \(S\) （ \(A\) 的）是乘性的且饱和的，充分必要条件是 \(A - S\) 是 \(A\) 的若干素理想的并。
\item
  设 S 是 A 的一个乘性子集，\(\tilde{S}\) 是由那些 \(x \in A\) 组成的集合，对这些 x 存在 \(y \in A\) 使 \(xy \in S\) 。证明 \(\tilde{S}\) 是包含 S 的最小饱和乘性子集，A --- \(\tilde{S}\) 是 A 中与 S 不相交的素理想的并，而且对每个 A-模 M，从 \(S^{-1}M\) 到 \(\tilde{S}^{-1}M\) 的典范映射是双射的。
\item
  设 S 与 T 是 A 的两个乘性子集。证明下面两个性质等价：(α) \(S \subset \tilde{T}\) ；(β) 对每个使 \(S^{-1}M = 0\) 的 A-模 M，有 \(T^{-1}M = 0\) 。（为看出 (β) 蕴含 (α)，考虑一个商 A-模 A/As，其中 \(s \in S\) 。）
\end{enumerate}

\begin{enumerate}
\def\labelenumi{\arabic{enumi}.}
\setcounter{enumi}{1}
\tightlist
\item
  设 A 是一个环，S 是 A 的乘性子集。证明由那些 \(x \in A\) （对某个 \(s \in S\) 满足 sx = 0）组成的集合 n 是 A 的理想。令 \(A_{1} = A/n\) ，并用 \(S_{1}\) 表示 S 在 \(A_{1}\) 中的典范像。证明 \(S_{1}\) 中没有元素是 \(\mathbf{A}_1\) 的零因子，并且典范同态 \(\mathbf{S}^{-1}\mathbf{A} \to \mathbf{S}_1^{-1}\mathbf{A}_1\) 是双射的。
\end{enumerate}

由此推出，为使 \(S^{-1}A\) 是有限生成的 A-模，充分必要条件是 \(S_{1}\) 的所有元素在 \(A_{1}\) 中可逆，这时 \(S^{-1}A\) 等同于 \(A_{1}\) 。

\begin{enumerate}
\def\labelenumi{\arabic{enumi}.}
\setcounter{enumi}{2}
\tightlist
\item
  设 \(S = Z^{*}\) （即 \(\{0\}\) 在 Z 中的补集），p 是一个整数 \textgreater1 ，M 是诸模 \(Z/p^{n}Z\) （ \(n \in N\) ）的直和所成的 Z-模。证明尽管 M 是忠实的 Z-模，却有 \(S^{-1}M = 0\) ，而且典范映射
\end{enumerate}

\[
\mathrm{S} ^ {- 1} \operatorname{End} _ {\mathbf {z}} (\mathrm{M}) \rightarrow \operatorname{End} _ {\mathrm{s} ^ {- 1} \mathbf {z}} (\mathrm{S} ^ {- 1} \mathrm{M})
\]

不是单射的。

\begin{enumerate}
\def\labelenumi{\arabic{enumi}.}
\setcounter{enumi}{3}
\item
  设 \(S = \mathbf{Z}^*\) ，\(M\) 是具有无穷基的自由 \(\mathbf{Z}\) -模。证明典范映射 \(S^{-1}\operatorname{End}_{\mathbf{Z}}(M) \to \operatorname{End}_{S^{-1}\mathbf{Z}}(S^{-1}M)\) 不是满射的。由此并借助习题 3，给出一个 \(\mathbf{Z}\) -模 \(M\) 的例子，使典范映射 \(S^{-1}\operatorname{End}_{\mathbf{Z}}(M) \to \operatorname{End}_{S^{-1}\mathbf{Z}}(S^{-1}M)\) 既非单射也非满射。
\item
  给出一个序列 \((\mathbf{P}_n)\) （ \(\mathbf{Z}\) 的子 \(\mathbf{Z}\) -模的序列）的例子，使得对 \(\mathbf{S} = \mathbf{Z}^*\) ，子模 \(\mathrm{S}^{-1}\left(\bigcap_{n} \mathrm{P}_{n}\right)\) 异于 \(\bigcap_{n} (\mathrm{S}^{-1}\mathrm{P}_{n})\) 。
\item
  给出一个 \(\mathbf{Z}\) -模 \(\mathbf{M}\) 的例子，使得对 \(\mathbf{S} = \mathbf{Z}^*\) ，\(\operatorname{Ann}(\mathbf{M}) = 0\) ，但 \(\operatorname{Ann}(\mathbf{S}^{-1}\mathbf{M}) = \mathbf{Q}\) （参见习题 3）。
\item
  设 S 是环 A 的一个乘性子集；对每个 A-模族 \((M_{t})_{t\in I}\) ，定义典范同态 \(S^{-1}\prod_{i\in I}M_{i}\to\prod_{i\in I}S^{-1}M_{i}\) （ \(S^{-1}A\) -模的）。给出一个使这个同态既非单射又非满射的例子（参见习题 3）。
\item
  设 A 是环，\((\mathrm{S}_{\lambda})_{\lambda \in L}\) 是 A 的乘性子集族，\(M = \bigoplus_{\lambda \in L} S_{\lambda}^{-1}A\) 。证明下面两个性质等价：( \(\alpha\) ) M 是忠实平坦 A-模；( \(\beta\) ) 对 A 的每个极大理想 m，存在 \(\lambda \in L\) 使 \(m \cap S_{\lambda} = \varnothing\) 。特别地，若 S 是 A 的乘性子集，则 A-模 \(S^{-1}A\) 只有当 S 的元素都可逆时才是忠实平坦的，这时 \(S^{-1}A\) 等同于 A。
\item
  设 A 是环，S 是 A 的乘性子集，q 是 A 的素理想。证明若 T 是 A - q 在 \(S^{-1}A\) 中的像，则环 \(T^{-1}(S^{-1}A)\) 同构于 \(S'^{-1}A\) ，其中 \(S'\) 是 A 中含于 q 且与 S 不相交的素理想的并的补集。
\item
  \begin{enumerate}
  \def\labelenumii{(\alph{enumii})}
  \tightlist
  \item
    设 K 是交换域，A 是多项式环 K{[}X, Y{]} ，\(p = (X)\) ，\(q = (Y)\) ，S 是 \(p \cup q\) 在 A 中的补集。证明 \(p + q\) 关于 S 的饱和异于 p 的饱和与 q 的饱和之和。
  \end{enumerate}
\end{enumerate}

\begin{enumerate}
\def\labelenumi{(\alph{enumi})}
\setcounter{enumi}{1}
\tightlist
\item
  设 K 是交换域，A 是多项式环 K{[}X, Y, Z{]} ，\(\mathfrak{p} = (\mathbf{X}) + (\mathbf{Z})\) ，\(\mathfrak{q} = (\mathbf{Y}) + (\mathbf{Z})\) ，S 是 \(p \cup q\) 在 A 中的补集。证明 pq 关于 S 的饱和异于 p 的饱和与 q 的饱和之积。
\end{enumerate}

\begin{enumerate}
\def\labelenumi{\arabic{enumi}.}
\setcounter{enumi}{10}
\item
  设 \(p\) 是素数；\(\mathbf{Z}\) 的哪些理想关于理想 \((p)\) 是饱和的？
\item
  在环 A 中，用 \(\mathfrak{r}(\mathfrak{a})\) 表示理想 \(\mathfrak{a}\) 的根。
\end{enumerate}

\begin{enumerate}
\def\labelenumi{(\alph{enumi})}
\tightlist
\item
  证明对 A 的三个理想 a、b、c，
\end{enumerate}

\[
\mathfrak {r} (\mathfrak {a} + \mathfrak {b c}) = \mathfrak {r} (\mathfrak {a} + (\mathfrak {b} \cap \mathfrak {c})) = \mathfrak {r} (\mathfrak {a} + \mathfrak {b}) \cap \mathfrak {r} (\mathfrak {a} + \mathfrak {c}).
\]

\begin{enumerate}
\def\labelenumi{(\alph{enumi})}
\setcounter{enumi}{1}
\tightlist
\item
  给出 A 的理想的一个无穷序列 \((\mathfrak{a}_n)\) 的例子，使得 \(\mathfrak{r}\left(\bigcap_{n} \mathfrak{a}_n\right) \neq \bigcap_{n} \mathfrak{r}(\mathfrak{a}_n)\) 。
\end{enumerate}

*(c) 给出一个主理想 \(\mathfrak{a}\) 的例子，使得 \(\mathfrak{r}(\mathfrak{a})\) 不是有限生成的，而且不存在整数 \(n > 0\) 使 \((\mathfrak{r}(\mathfrak{a}))^n \subset \mathfrak{a}\) 。（考虑一个高为 1 的非离散赋值的环。）*

13. 设 A 是一个环.\par

\begin{enumerate}
\def\labelenumi{(\alph{enumi})}
\item
  通过考虑多项式 \(f\) 在环 \((\mathbf{A} / \mathfrak{p})[\mathbf{X}]\) 中的像（其中 \(\mathfrak{p}\) 取遍 \(\mathbf{A}\) 的素理想集合），给出《代数》第 VIII 章 § 6 习题 6 (b) 的另一个证明。
\item
  设 \(N\) 是一个 \(r\) 阶方阵（在 \(A\) 上）；为使 \(N\) 是幂零的，充分必要条件是它的特征多项式的每个系数（首项系数除外）都是幂零的。（为看出条件必要，用与 (a) 相同的方法；为看出条件充分，用 Cayley-Hamilton 定理。）
\item
  证明对每对正整数 \(k, r\) ，存在一个最小正整数 \(m(k, r)\) （它不依赖于环 \(\mathbf{A}\) ），使得关系 \(N^k = 0\) 蕴含 \((\operatorname{Tr}(N))^{m(k,r)} = 0\) 。（考虑环 \(\Lambda = \mathbf{Z}[T_{ij}]\) ，其中诸 \(T_{ij}\) 是 \(r^2\) 个未定元，以及矩阵 \(X = (T_{ij})\) （ \(\Lambda\) 上的）与元素 \(t = \sum_{i=1}^{n} \mathrm{T}_{ii} = \mathrm{Tr}(X)\) （ \(\Lambda\) 的）；若 \(\mathfrak{a}_k\) 是 \(\Lambda\) 的由 \(X^k\) 的元素生成的理想，则借助 (b) 证明存在整数 \(m\) 使 \(t^m \in \mathfrak{a}_k\) 。）
\item
  证明 \(m(2, 2) = 4\) 。
\end{enumerate}

¶ 14. (a) 设 A 是左 Noether 环（不必交换），a 是 A 的一个左理想。假设存在左理想 b ≠ 0 使 a ∩ b = 0。

证明每个元素 \(a \in \mathfrak{a}\) 都是右零因子。（设 \(b \neq 0\) 是 b 的一个元素，令 \(a_{n} = Ab + Aba + \cdots + Aba^{n}\) ；考虑使 \(a_{n} = a_{n+1}\) 的最小整数 n。

\begin{enumerate}
\def\labelenumi{(\alph{enumi})}
\setcounter{enumi}{1}
\tightlist
\item
  设 E 是域 K 上的无穷维向量空间，A 是环 \(\operatorname{End}_{\mathbb{K}}(\operatorname{E})\) 。给出 A 的两个元素 u、v 的例子，使 u 与 v 都不是右零因子且 \(Au \cap Av = 0\) 。
\end{enumerate}

¶ 15. 设 X 是完全正则拓扑空间，βX 是它的 Stone-Čech 紧化（《一般拓扑学》第 IX 章 § 1 习题 7）。对所有 \(x \in \mathbf{X}\) ，用 \(\mathfrak{J}(x)\) 表示环 \(\mathcal{C}(\mathbf{X}; \mathbf{R})\) 中由那些 \(f \in \mathcal{C}(\mathbf{X}; \mathbf{R})\) 组成的极大理想（对这些 f，\(x\) 属于 X 的子集 \(\bar{f}^1(0)\) 在 X 中的闭包）（《一般拓扑学》第 X 章 § 4 习题 15）；用 \(\mathfrak{U}(x)\) 表示由那些 \(f \in \mathcal{C}(\mathbf{X}; \mathbf{R})\) 组成的理想（对这些 f，\(\bar{f}^1(0)\) 是 \(x\) 在 X 中的某个邻域在 X 上的迹）；\(\mathfrak{U}(x)\) 等于它的根。

理想 \(\mathfrak{a}\) （ \(\mathfrak{C}(\mathbf{X};\mathbf{R})\) 的）称为孤立的（相应地，绝对孤立的），如果对所有 \(f\geqslant 0\) （ \(\mathfrak{a}\) 中的），每个 \(g\) （满足 \(0\leqslant g\leqslant f\) （相应地，\(|g|\leqslant f\) ）的）都属于 \(\mathfrak{a}\) （参见《代数》第 IV 章 §1 习题 4）；这时 \(\mathcal{C}(\mathbf{X};\mathbf{R}) / \mathfrak{a}\) 具有一个与其环结构相容的序结构（相应地，格序结构），并且其 \(\geqslant 0\) 的元素是模 \(\mathfrak{a}\) 的类：即 \(\geqslant 0\) 的元素（ \(\mathcal{C}(\mathbf{X};\mathbf{R})\) 中的）的类（同上）。

\begin{enumerate}
\def\labelenumi{(\alph{enumi})}
\item
  证明理想 \(\mathfrak{a}\) （ \(\mathcal{C}(\mathbf{X};\mathbf{R})\) 的、等于其自身的根的）是绝对孤立的（若 \(f\in \mathfrak{a}\) 且 \(|g|\leqslant |f|\) ，考虑当 \(f(x) = 0\) 时等于 0、否则等于 \((g(x))^2 /f(x)\) （此时 \(f(x)\neq 0\) ）的函数）。在这种情形，为使 \(\mathcal{C}(\mathbf{X};\mathbf{R}) / \mathfrak{a}\) 是全序的，充分必要条件是 \(\mathfrak{a}\) 为素理想；这时 \(\mathcal{C}(\mathbf{X};\mathbf{R}) / \mathfrak{a}\) 的素理想集按包含关系是全序的。
\item
  为使孤立理想 \(\mathfrak{a}\) （ \(\mathcal{C}(\dot{\mathbf{X}};\mathbf{R})\) 的）满足 \(\mathcal{C}(\mathbf{X};\mathbf{R}) / \mathfrak{a}\) 是全序的，必要的是 \(\mathfrak{a}\) 含于唯一的极大理想 \(\mathfrak{J}(x)\) 之中，且这时 \(\mathfrak{a}\) 必然包含 \(\mathfrak{U}(x)\) ；充分的是 \(\mathfrak{a}\) 包含一个素理想，这时 \(\mathfrak{a}\) 是绝对孤立的（注意这时关系 \(fg = 0\) 蕴含 \(f\in \mathfrak{a}\) 或 \(g\in \mathfrak{a}\) ）。
\item
  取 \(\mathbf{X} = \mathbf{R}\) ，\(x = 0\) ；证明由函数 \(t \mapsto |t|\) 生成的主理想含于 \(\mathfrak{J}(0)\) 且包含 \(\mathfrak{U}(0)\) ，但不是孤立的；由恒等映射 \(t \mapsto t\) 生成的主理想含于 \(\mathfrak{J}(0)\) ，包含 \(\mathfrak{U}(0)\) ，并且是孤立的但不是绝对孤立的。
\item
  证明若 \(\mathfrak{p}\) 与 \(\mathfrak{q}\) 是 \(\mathcal{C}(\mathbf{X};\mathbf{R})\) 中两个素理想，则 \(\mathfrak{pq} = \mathfrak{p} \cap \mathfrak{q}\) （特别地，\(\mathfrak{p}^2 = \mathfrak{p}\) ），并且 \(\mathfrak{p} + \mathfrak{q}\) 是素的或等于 \(\mathcal{C}(\mathbf{X};\mathbf{R})\) 。（若 \(\mathfrak{p} + \mathfrak{q} \neq \mathcal{C}(\mathbf{X};\mathbf{R})\) ，注意 \(\mathfrak{p} + \mathfrak{q}\) 是绝对孤立的且等于它的根，然后利用 (a)。）
\item
  为使 \(\mathfrak{U}(x)\) 是素的，充分必要条件是：对每个函数 \(f\in\mathcal{C}(\mathbf{X};\mathbf{R})\) ，在 \(\beta\mathbf{X}\) 中，两个集合 \(\bar{f}^{-1}\left([0,+\infty[),\bar{f}^{-1}([- \infty,0])\right)\) 之一的闭包是 \(x\) 的邻域。由此推出一个使 \(\mathfrak{U}(x)\) 是素的但非极大的例子（考虑与一个非平凡超滤子相伴的拓扑空间；参见《一般拓扑学》第 I 章 § 6 第 5 节）。
\item
  设 \(x \in \mathbf{X}\) 。为使 \(\mathfrak{U}(x) = \mathfrak{J}(x)\) ，充分必要条件是 \(x\) 在 \(\mathbf{X}\) 中的邻域的每个可数交仍是 \(x\) 在 \(\mathbf{X}\) 中的邻域。（参见《一般拓扑学》第 I 章 § 7 习题 7。）
\item
  证明 \(\mathcal{C}(\mathbf{X};\mathbf{R})\) 典范地等同于 \(\mathcal{C}^{\infty}(\mathbf{X};\mathbf{R})\) 的全分式环。
\end{enumerate}

\begin{enumerate}
\def\labelenumi{\arabic{enumi}.}
\setcounter{enumi}{15}
\tightlist
\item
  设 A 是一个拓扑环（不必交换）。A 上的拓扑称为（左）线性的，如果 A 的开左理想构成 0 的一个基本邻域系。若如此，则 A 的开左理想所成的集合 F 满足下列条件：
\end{enumerate}

\begin{enumerate}
\def\labelenumi{(\arabic{enumi})}
\item
  A 的每个包含某个理想 \(m \in \mathcal{F}\) 的左理想都属于 \(\mathcal{F}\) 。
\item
  \(\mathcal{F}\) 中理想的每个有限交属于 \(\mathcal{F}\) 。
\item
  若 \(\mathfrak{m} \in \mathcal{F}\) 且 \(a \in \mathbf{A}\) ，则左理想 \(\mathfrak{n}\) （由那些 \(x \in \mathbf{A}\) ——满足 \(xa \in \mathfrak{m}\) ——组成）属于 \(\mathcal{F}\) 。
\end{enumerate}

反之，若环 A 的左理想的一个集合 F 满足这三个条件，则存在唯一的与 A 上的环结构相容且使 F 为 A 的开左理想集合的拓扑。这样的左理想集合 F 称为（左）拓扑化集。

\begin{enumerate}
\def\labelenumi{\arabic{enumi}.}
\setcounter{enumi}{16}
\tightlist
\item
  (*) 设 A 是一个环（不必交换），\(\mathcal{F}\) 是 A 的左理想的一个拓扑化集（习题 16）。
\end{enumerate}

\begin{enumerate}
\def\labelenumi{(\alph{enumi})}
\item
  （左）A-模 M 称为 F-可忽略的，如果 M 的每个元素的零化子都属于 F。若 M 是 F-可忽略的，则它的每个子模和每个商模也是 F-可忽略的。若 M 与 N 是两个 F-可忽略的 A-模，则 \(M \oplus N\) 是 F-可忽略的。为使 \(A_{s}/m\) 是 F-可忽略的，充分必要条件是 \(m \in F\) 。若 \((0) \in F\) ，则每个 A-模都是 F-可忽略的；若 F 只由单个左理想 \(A_{s}\) 组成，则每个 F-可忽略的 A-模都化为 0。
\item
  对每个 A-模 M，存在 M 的一个最大的 F-可忽略子模，记作 FM。对每个 A-模同态 \(u: M \to N\) ，有 \(u(\mathcal{F}M) \subset \mathcal{F}N\) ；用 Fu 表示从 FM 到 FN 的、图象与 \(u|FM\) 的图象重合的那个映射。若
\end{enumerate}

\[
0 \longrightarrow \mathrm{M} \stackrel {u} {\longrightarrow} \mathrm{N} \stackrel {v} {\longrightarrow} \mathrm{P}
\]

是 A-模的一个正合列，则序列

\[
0 \longrightarrow \mathcal {F} M \stackrel {\mathcal {F} _ {u}} {\longrightarrow} \mathcal {F} N \stackrel {\mathcal {F} _ {v}} {\longrightarrow} \mathcal {F} P
\]

是正合的。证明 \(FA_{s}\) 是 A 的一个双边理想。

\begin{enumerate}
\def\labelenumi{(\alph{enumi})}
\setcounter{enumi}{2}
\tightlist
\item
  设 \(\mathfrak{m}\) 、\(\mathfrak{n}\) 是 \(\mathcal{F}\) 的两个元素且 \(\mathfrak{m} \subset \mathfrak{n}\) ；对每个左 A-模 \(\mathbf{M}\) ，典范单射 \(j \colon \mathfrak{m} \to \mathfrak{n}\) 定义一个交换群同态
\end{enumerate}

\[
u _ {m, n} = \operatorname{Hom} _ {A} (j, 1 _ {M}): \operatorname{Hom} _ {A} (n, M) \rightarrow \operatorname{Hom} _ {A} (m, M).
\]

若 \(\mathcal{F}\) 按关系 \(\supset\) 排序，则诸 \(u_{m,n}\) 定义一个交换群的正向系，其正向极限记作 \(\mathrm{M}_{(\mathcal{F})}\) 。诸映射 \(u_{m,A_s}\) 构成一个正向系，其正向极限是一个从 \(\mathrm{M} = \mathrm{Hom}_{\mathrm{A}}(\mathrm{A}_s,\mathrm{M})\) 到 \(\mathrm{M}_{(\mathcal{F})}\) 的映射（称为典范映射）。类似地，若 \(v: \mathrm{M} \to \mathrm{N}\) 是一个左 A-模同态，则诸 \(v_m = \mathrm{Hom}_{\mathrm{A}}(1,v): \mathrm{Hom}_{\mathrm{A}}(\mathfrak{m},\mathrm{M}) \to \mathrm{Hom}_{\mathrm{A}}(\mathfrak{m},\mathrm{N})\) 构成一个正向系，其正向极限是一个交换群同态

\[
v _ {(\mathcal {F})} \colon \mathrm{M} _ {(\mathcal {F})} \rightarrow \mathrm{N} _ {(\mathcal {F})}.
\]

若 \(0 \to \mathbf{M} \xrightarrow{v} \mathbf{N} \xrightarrow{w} \mathbf{P}\) 是 A-模的一个正合列，则序列

\[
0 \longrightarrow \mathrm{M} _ {(\mathcal {F})} \stackrel {v (\mathcal {F})} {\longrightarrow} \mathrm{N} _ {(\mathcal {F})} \stackrel {w (\mathcal {F})} {\longrightarrow} \mathrm{P} _ {(\mathcal {F})}
\]

是正合的。

\begin{enumerate}
\def\labelenumi{\arabic{enumi}.}
\setcounter{enumi}{17}
\tightlist
\item
  \begin{enumerate}
  \def\labelenumii{(\alph{enumii})}
  \tightlist
  \item
    设 A 是一个环，F 与 G 是 A 的左理想的两个拓扑化集（习题 16）。用 F.G 表示 A 的满足如下条件的左理想 l 组成的集合：存在左理想 n ∈ G 使 n/l 是 F-可忽略的。证明 F.G 是拓扑化的；为使一个 A-模 M 是 F.G-可忽略的，充分必要条件是存在 M 的一个 F-可忽略子模 M' 使 M/M' 是 G-可忽略的。证明 A 的左理想的拓扑化集全体上的合成规律 (F,G) → F.G 是结合的。
  \end{enumerate}
\end{enumerate}

称 \(\mathcal{F}\) 是幂等的，如果 \(\mathcal{F}.\mathcal{F} = \mathcal{F}\) 。

\begin{enumerate}
\def\labelenumi{(\alph{enumi})}
\setcounter{enumi}{1}
\item
  证明 \(\mathcal{F}.\mathcal{G}\) 包含所有形如 \(\mathfrak{m}.\mathfrak{n}\) 的理想，其中 \(\mathfrak{m}\in\mathcal{F}\) ，\(\mathfrak{n}\in\mathcal{G}\) 。
\item
  设 K 是一个交换域，B 是系数在 K 中、无穷多个未定元 \((\mathbf{X}_{i})\) ( \(i \geqslant 0\) ) 上的多项式环。设 b 是 B 的由诸元素 \(X_{i}X_{j}\) ( \(i \neq j\) ) 生成的理想；设 A 是环 B/b，\(\xi_{i}\) 是 \(X_{i}\) 模 b 的类，m 是 A 的由诸 \(\xi_{i}\) 生成的理想；取 F 为包含 m 的某个幂的那些理想组成的集合。证明 F 是拓扑化的且包含 F 中任意两个理想之积，但不是幂等的。
\item
  假设 A 是交换的，F 是 A 的理想的一个拓扑化集，使得每个理想 \(m \in F\) 都包含一个有限生成理想 \(n \in F\) 。证明若 F 中两个理想之积属于 F，则 F 是幂等的。
\end{enumerate}

¶ 19. 设 F 是环 A 的左理想的一个集合；假设 F 是拓扑化的且幂等的（习题 16 与 18）。

\begin{enumerate}
\def\labelenumi{(\alph{enumi})}
\tightlist
\item
  证明若 \(\mathfrak{m} \in \mathcal{F}\) ，\(\mathfrak{n} \in \mathcal{F}\) 且 \(u \in \operatorname{Hom}_{\mathbf{A}}(\mathfrak{n}, \mathbf{A}_s)\) ，则 \(\bar{u}^1(\mathfrak{m}) \in \mathcal{F}\) （考虑正合列 \(0 \to \mathfrak{n} / \bar{u}^1(\mathfrak{m}) \to \mathbf{A}_s / \bar{u}^1(\mathfrak{m}) \to \mathbf{A}_s / \mathfrak{n} \to 0\) ）。对每个 A-模 M 以及所有 \(v \in \operatorname{Hom}_{\mathbf{A}}(\mathfrak{m}, \mathbf{M})\) ，用 \(u.v\) 表示 \(\mathbf{M}_{(\mathcal{F})}\) （习题 17 (c)）中复合同态 \(\bar{u}^1(\mathfrak{m}) \xrightarrow{u} \mathfrak{m} \xrightarrow{v} \mathbf{M}\) 的典范像。证明映射 \((u,v) \mapsto u.v\) 是 Z-双线性的；与之对应有一个 Z-线性映射
\end{enumerate}

\[
\psi_ {m, n} \colon \operatorname{Hom} _ {A} (n, A _ {s}) \otimes_ {\mathbf {Z}} \operatorname{Hom} _ {A} (m, M) \rightarrow M _ {(\mathcal {F})}.
\]

这些映射构成一个正向系，因而它们的正向极限典范地对应于一个 Z-双线性映射

\[
\psi \colon A _ {(\mathcal {F})} \times M _ {(\mathcal {F})} \rightarrow M _ {(\mathcal {F})}.
\]

证明若 \(M = A_{s}\) ，则 \(\psi\) 是一个内合成规律，它使 \(\mathbf{A}_{(\mathcal{F})}\) 成为一个环，而且典范映射 \(A \to A_{(\mathcal{F})}\) 对这个结构是环同态。对任意左 A-模 M，\(\psi\) 是一个外规律，它在 \(\mathbf{M}_{(\mathcal{F})}\) 上定义一个左 \(\mathbf{A}_{(\mathcal{F})}\) -模结构；对这个结构（以及典范映射 \(A \to A_{(\mathcal{F})}\) ），典范映射 \(M \to M_{(\mathcal{F})}\) 是 A-模同态。

\begin{enumerate}
\def\labelenumi{(\alph{enumi})}
\setcounter{enumi}{1}
\tightlist
\item
  若 \(i: \mathbf{M} \to \mathbf{M}_{(\mathcal{F})}\) 是典范同态，证明
\end{enumerate}

\[
\operatorname{Ker} (i) = \mathcal {F M}
\]

并且 A-模 Coker \((i) = \mathbf{M}_{(\mathcal{F})} / i(\mathbf{M})\) 是 \(\mathcal{F}\) -可忽略的。

\begin{enumerate}
\def\labelenumi{(\alph{enumi})}
\setcounter{enumi}{2}
\tightlist
\item
  对每个左 A-模 M，用 \(M_{F}\) 表示左 A-模 \((\mathrm{M}/\mathcal{F}\mathrm{M})_{(\mathcal{F})}\) ，用 \(j_{M}\) 表示复合的典范映射
\end{enumerate}

\[
\mathbf {M} \rightarrow \mathbf {M} / \mathcal {F} \mathbf {M} \rightarrow \mathbf {M} _ {\mathcal {F}};
\]

于是 \(\operatorname{Ker}(j_{\mathbf{M}}) = \mathcal{F}\mathbf{M}\) ，并且 A-模 \(\operatorname{Coker}(j_{\mathbf{M}}) = \mathrm{M}_{\mathcal{F}} / j_{\mathbf{M}}(\mathbf{M})\) 是 \(\mathcal{F}\) -可忽略的。

设 \(u: \mathrm{P} \to \mathrm{M}\) ，\(h: \mathrm{P} \to \mathrm{Q}\) 是 A-模同态，使得 \(\operatorname{Ker}(h)\) 与 \(\operatorname{Coker}(h)\) 都是 \(\mathcal{F}\) -可忽略的；证明存在唯一的 A-同态 \(v: \mathrm{Q} \to \mathrm{M}_{\mathcal{F}}\) 使下图交换

\[
\begin{array}{c} \mathrm{P} \xrightarrow {u} \mathrm{M} \\ h \Big \downarrow \qquad \qquad \qquad \qquad \qquad \qquad \qquad \Big \downarrow j _ {\mathrm{M}} \\ \mathrm{Q} \xrightarrow [ v ]{} \mathrm{M} _ {\mathcal {F}} \end{array}
\]

（先化为 h 是单射的情形；然后，把 P 等同于 Q 的一个子模，注意对 \(x \in Q\) ，存在 \(m \in F\) 使 \(mx \in P\) ，并把 x 映为复合同态在 \(M_{F}\) 中的典范像

\[
\mathfrak {m} \xrightarrow {x} \mathfrak {m} x \to \mathrm{P} \xrightarrow {u} \mathrm{M} \to \mathrm{M} / \mathcal {F} \mathrm{M.})
\]

\begin{enumerate}
\def\labelenumi{(\alph{enumi})}
\setcounter{enumi}{3}
\tightlist
\item
  对每个 A-模同态 \(u: M \rightarrow N\) ，证明存在唯一的同态 \(u_{F}: M_{F} \rightarrow N_{F}\) 使下图交换
\end{enumerate}

\[
\begin{array}{c c c} \mathrm{M} & \stackrel {{u}} {{\longrightarrow}} \mathrm{N} \\ j _ {\mathrm{M}} \Big \downarrow & & \Big \downarrow j _ {\mathrm{N}} \\ \mathrm{M} _ {\mathcal {F}} & \stackrel {{u _ {\mathcal {F}}}} {{\longrightarrow}} \mathrm{N} _ {\mathcal {F}} \end{array}
\]

（利用 (a)、(b) 与 (c)）。证明若

\[
0 \longrightarrow \mathrm{M} \stackrel {u} {\longrightarrow} \mathrm{N} \stackrel {v} {\longrightarrow} \mathrm{P}
\]

是 A-模的一个正合列，则序列

\[
0 \longrightarrow \mathrm{M} _ {\mathcal {F}} \stackrel {u \mathcal {F}} {\longrightarrow} \mathrm{N} _ {\mathcal {F}} \stackrel {v \mathcal {F}} {\longrightarrow} \mathrm{P} _ {\mathcal {F}}
\]

是正合的（注意映射 \(\mathbf{M} / \mathcal{F}\mathbf{M}\to \mathrm{N} / \mathcal{F}\mathrm{N}\) ——由 \(u\) 取商得到——是单射的）。

证明映射 \(j_{\mathbf{M}_{\mathcal{F}}}\) 与 \((j_{\mathbf{M}})_{\mathcal{F}}\) 重合且都是同构（因而 \(\mathcal{FM}_{\mathcal{F}} = 0\) ）。若 \(\mathbf{N}'\) 是 \(\mathbf{M}_{\mathcal{F}}\) 的一个子 A-模，则 \((\overline{j_{\mathbf{M}}^{-1}(\mathbf{N}')})_{\mathcal{F}}\) 典范地等同于 \(\mathbf{N}_{\mathcal{F}}'\) （注意，若记 \(\mathbf{N} = j_{\mathbf{M}}^{-1}(\mathbf{N}')\) ，则 \(\mathbf{N}' / j_{\mathbf{M}}(\mathbf{N})\) 是 \(\mathcal{F}\) -可忽略的）。

\begin{enumerate}
\def\labelenumi{(\alph{enumi})}
\setcounter{enumi}{4}
\tightlist
\item
  设 M 是一个左 A-模，\(\phi: \mathbf{A} \times \mathbf{M} \to \mathbf{M}\) 是 Z-双线性映射 \((a, m) \mapsto am\) 。证明存在唯一的 Z-双线性映射 \(\phi_{\mathcal{F}}: \mathbf{A}_{\mathcal{F}} \times \mathbf{M}_{\mathcal{F}} \to \mathbf{M}_{\mathcal{F}}\) 使下图交换
\end{enumerate}

\[
\begin{array}{c} \mathbf {A} \times \mathbf {M} \xrightarrow {\phi} \mathbf {M} \\ j _ {\mathbf {A}} \times j _ {\mathbf {M}} \Big \downarrow \qquad \qquad \qquad \qquad \Big \downarrow j _ {\mathbf {M}} \\ \mathbf {A} _ {\mathcal {F}} \times \mathbf {M} _ {\mathcal {F}} \xrightarrow [ \phi_ {\mathcal {F}} ]{} \mathbf {M} _ {\mathcal {F}} \end{array}
\]

同样在 \(A_{F}\) 上定义一个环结构，并在 \(M_{F}\) 上定义一个左 \(A_{F}\) -模结构。若 \(u: M \to N\) 是 A-模同态，则 \(u_{F}: M_{F} \to N_{F}\) 是 \(A_{F}\) -模同态。为使对每个 A-模 M，映射 \(u \mapsto u_{F}\) （从 \(\operatorname{Hom}_{\mathbf{A}}(\mathbf{M}, \mathbf{N})\) 到 \(\operatorname{Hom}_{\mathbf{A}_{F}}(\mathbf{M}_{F}, \mathbf{N}_{F})\) 的）是双射的，充分必要条件是 \(j_{N}\) 是双射的。

\begin{enumerate}
\def\labelenumi{(\alph{enumi})}
\setcounter{enumi}{5}
\item
  取 A 为交换域 K 上的多项式环 K{[}X, Y{]} ，F 为 A 的包含理想 m = (X) + (Y) 的某个幂的理想组成的（拓扑化且幂等的）集合。若 M = A/AX 且 u: A → M 是典范映射，证明 \(u_{F}: A_{F} \to M_{F}\) 不是满射的。（证明 \(A_{F}\) 等同于 A，而 \(M_{F}\) 等同于 K{[} \(\overline{Y}\) , 1/ \(\overline{Y}\) {]} ，其中 \(\overline{Y}\) 是 Y 在 M 中的类，于是 X \(\overline{Y}\) = 0 且 X(1/ \(\overline{Y}\) ) = 0。）
\item
  若令 \(\mathcal{F}^1\mathrm{M} = \mathrm{M}_{\mathcal{F}} / j_{\mathrm{M}}(\mathrm{M})\) ，证明对 A-模的每个正合列 \(0\to \mathrm{M}\to \mathrm{N}\to \mathrm{P}\to 0\) ，有一个正合列
\end{enumerate}

\[
0 \rightarrow \mathcal {F} M \rightarrow \mathcal {F} N \rightarrow \mathcal {F} P \rightarrow \mathcal {F} ^ {1} M \rightarrow \mathcal {F} ^ {1} N \rightarrow \mathcal {F} ^ {1} P
\]

（利用蛇形图）。

\begin{enumerate}
\def\labelenumi{(\alph{enumi})}
\setcounter{enumi}{7}
\tightlist
\item
  设 \(\mathcal{F}'\) 是由左理想 \(l'\) （ \(\mathbf{A}_{\mathcal{F}}\) 中满足 \((\mathbf{A}_{\mathcal{F}})_s / l'\) 是 \(\mathcal{F}\) -可忽略 A-模的）组成的集合：证明 \(\mathcal{F}'\) 是 \(\mathbf{A}_{\mathcal{F}}\) 的包含 \(j_{\mathbf{A}}(\mathfrak{m})\) （其中 \(\mathfrak{m} \in \mathcal{F}\) ）的左理想组成的集合（注意，利用 (d) 的正合列，有 \(\mathfrak{m}_{\mathcal{F}} = \mathbf{A}_{\mathcal{F}}\) （对所有 \(\mathfrak{m} \in \mathcal{F}\) ）；因而 \(A_{\mathcal{F}}/j_{A}(m)\) 是 F-可忽略的）。对每个 \(A_{F}\) -模 \(M'\) ，有 \(FM' = F'M'\) （ \(M'\) 借助 \(j_{A}\) 被看作 A-模），\(F'\) 是拓扑化的且幂等的，并且 \(M'_{F'}\) 典范地等同于 \(M'_{F}\) 。
\end{enumerate}

\begin{enumerate}
\def\labelenumi{\arabic{enumi}.}
\setcounter{enumi}{19}
\tightlist
\item
  设 \(\mathcal{F}\) 是环 A 的左理想的一个拓扑化且幂等的集合。
\end{enumerate}

\begin{enumerate}
\def\labelenumi{(\alph{enumi})}
\item
  假设 \(\mathcal{F}\) 的每个理想都包含一个属于 \(\mathcal{F}\) 的有限生成理想。设 M 是一个 A-模，\((\mathbf{N}_{\iota})_{\iota \in I}\) 是 M 的子模的一个右有向族，其并为 M。证明 \(\mathbf{M}_{\mathcal{F}}\) 是它的子 \(\mathbf{A}_{\mathcal{F}}\) -模 \((\mathbf{N}_{\iota})_{\mathcal{F}}\) 的并。特别地，对 A-模的每个族 \((\mathbf{M}_{\lambda})_{\lambda \in L}\) ，\(\left( \bigoplus_{\lambda \in L} \mathbf{M}_{\lambda} \right)_{\mathcal{F}}\) 典范地同构于 \(\bigoplus_{\lambda \in L} (\mathbf{M}_{\lambda})_{\mathcal{F}}\) 。
\item
  假设条件 (a) 成立，并且对每个满射的 A-模同态 \(u \colon \mathbf{M} \to \mathbf{N}\) ，\(u_{\mathcal{F}} \colon \mathbf{M}_{\mathcal{F}} \to \mathbf{N}_{\mathcal{F}}\) 是满射的。这时，典范映射 \(\mathbf{A}_{\mathcal{F}} \otimes_{\mathbf{A}} \mathbf{M} \to \mathbf{M}_{\mathcal{F}}\) （由 \(\mathbf{M}_{\mathcal{F}}\) 作为 \(\mathbf{A}_{\mathcal{F}}\) -模的外规律所定义）是一个 \(\mathbf{A}_{\mathcal{F}}\) -模同构（化为 \(\mathbf{M}\) 是自由的情形）；环 \(\mathbf{A}_{\mathcal{F}}\) 是平坦的左 A-模，而且左 \(\mathbf{A}_{\mathcal{F}}\) -模与左 A-模 \(\mathbf{M}\) ——使映射 \(j_{\mathbf{M}} \colon \mathbf{M} \to \mathbf{M}_{\mathcal{F}}\) 为双射的——重合。
\item
  假设 A 是环的一个无穷积 \(\biguplus_{i\in I}A_{i}\) ；用 \(e_{i}\) 表示 \(A_{i}\) 的单位元素，取 F 为 A 的包含双边理想 \(a=\bigoplus_{i\in I}A_{i}\) 的左理想组成的集合（诸 \(A_{i}\) 典范地等同于 A 的理想）。证明对每个左 A-模 M，\(M_{F}\) 等同于积 \(\prod_{i\in I}(e_{i}M)\) ，特别地 \(A_{F}=A\) 。由此推出，对每个满射的 A-同态 \(u: M\to N\) ，\(u_{F}: M_{F}\to N_{F}\) 是满射的，但给出使 \(j_{M}\) 不是双射的 A-模 M 的例子。
\item
  假设条件 (a) 成立并且 A 是交换的。证明若 M 是 F-可忽略的，则 \(\mathbf{M}_{(\mathcal{F})}=0\) 。
\end{enumerate}

\begin{enumerate}
\def\labelenumi{\arabic{enumi}.}
\setcounter{enumi}{20}
\tightlist
\item
  设 A 是一个交换环，\(\mathcal{F}\) 是 A 的理想的一个拓扑化且幂等的集合。
\end{enumerate}

\begin{enumerate}
\def\labelenumi{(\alph{enumi})}
\tightlist
\item
  这时环 \(\mathbf{A}_{\mathcal{F}}\) 是交换的（注意，若 \(\mathfrak{m} \in \mathcal{F}\) 且
\end{enumerate}

\[
u \in \operatorname{Hom} _ {\mathbf {A}} (\mathfrak {m}, \mathbf {A}),
\]

则 \(u(\mathfrak{m}^2) \subset \mathfrak{m}\) ；若 \(v\) 是 \(\operatorname{Hom}_{\mathbf{A}}(\mathfrak{m}, \mathbf{A})\) 的另一个元素，证明

\[
v (u (x y)) = u (v (x y))
\]

对 \(x \in \mathfrak{m}, y \in \mathfrak{m}\) 成立）。

\begin{enumerate}
\def\labelenumi{(\alph{enumi})}
\setcounter{enumi}{1}
\item
  设 \(\mathcal{F}'\) 是理想 \(l'\) （ \(\mathbf{A}_{\mathcal{F}}\) 的，满足 \(\mathbf{A}_{\mathcal{F}} / l'\) 是 \(\mathcal{F}\) -可忽略的）组成的族（习题 19 (h)）。证明映射 \(\mathfrak{p} \mapsto \mathfrak{p}_{\mathcal{F}}\) 是理想 \(\mathfrak{p}\) （ \(\mathbf{A}\) 的、不属于 \(\mathcal{F}\) 的）全体到 \(\mathbf{A}_{\mathcal{F}}\) 的不属于 \(\mathcal{F}'\) 的素理想全体上的一个双射。（先证明若 \(\mathbf{A}\) 是整环，则 \(\Lambda_{\mathcal{F}}\) 也是整环，作与 (a) 中相同的注记；然后利用习题 19 (d) 的正合列证明 \(\mathfrak{p}_{\mathcal{F}}\) 是素的。若 \(\mathfrak{p}'\) 是 \(A_{\mathcal{F}}\) 的一个不属于 \(\mathcal{F}'\) 的素理想且 \(\mathfrak{p} = j_{\mathrm{A}}^{-1}(\mathfrak{p}')\) ，则回忆 \(\mathfrak{p}' = \mathfrak{p}_{\mathcal{F}}\) 。
\item
  设 B 是一个 A-代数，G 是 B 的满足 \(B_{s}/!\) 是 F-可忽略的左理想 l 组成的集合。证明 G 是 B 的包含一个形如 B.m 的理想（其中 \(m \in F\) ）的左理想组成的集合；由此推出 G 是拓扑化的且幂等的，并且对每个 B-模 N，有 FN = GN 且 \(N_{F}\) 典范地等同于 \(N_{G}\) 。
\item
  设 S 是 A 的一个乘性子集，F 是 A 的与 S 相交的理想组成的集合。证明 F 是拓扑化的且幂等的，并且对每个 A-模 M，\(M_{(F)}\) 与 \(M_{F}\) 都典范地等同于 \(S^{-1}M\) （注意 F 的每个理想都包含一个主理想 As，其中 \(s \in S\) ）。
\end{enumerate}

¶ 22. 设 A 是一个环（不必交换），S 是 A 的一个乘性子集。

\begin{enumerate}
\def\labelenumi{(\alph{enumi})}
\item
  设 \(\mathcal{F}\) 是由左理想 \(\mathfrak{l}\) （ \(\mathbf{A}\) 中满足如下条件的：对所有 \(a\in \mathbf{A}\) ，存在 \(s\in \mathbf{S}\) 使 \(sa\in \mathfrak{l}\) ）组成的集合；这特别蕴含 \(\mathfrak{l}\cap \mathbf{S}\neq \varnothing\) 。证明 \(\mathcal{F}\) 是拓扑化的且幂等的。
\item
  设 B 是一个环，\(\phi: A \to B\) 是一个环同态。若序偶 \((B, \phi)\) 满足下列条件，则称之为 \(A\) 的分母在 \(S\) 中的左分式环：
\end{enumerate}

\begin{enumerate}
\def\labelenumi{(\Roman{enumi})}
\item
  若 \(\phi(a) = 0\) ，则存在 \(s \in S\) 使 \(sa = 0\) 。
\item
  若 \(s \in S\) ，则 \(\phi(s)\) 在 \(B\) 中可逆。
\item
  B 的每个元素都形如 \((\phi(s))^{-1}\phi(a)\) ，其中 \(a \in \mathbf{A}\) 。
\end{enumerate}

证明下列性质等价：

(α) 环 A 有一个分母在 S 中的左分式环。

(β) 下列条件成立：

\((\beta_{1})\) 对所有 \(s \in S\) ，\(a \in A\) ，存在 \(t \in S\) 与 \(b \in A\) 使 \(ta = bs\) 。

\((\beta_{2})\) 若 \(a \in \mathbf{A}\) 与 \(s \in \mathbf{S}\) 满足 \(as = 0\) ，则存在 \(t \in \mathbf{S}\) 使 \(ta = 0\) 。

(γ) S 的诸元素在 \(A_{F}\) 中的典范像是可逆的。

此外这些性质蕴含下列性质：

(δ) 主理想 As（其中 \(s \in S\) ）属于 F，并且每个理想 \(m \in F\) 都包含这些主理想之一。

(ε) 每个 \(s \in S\) 的左零化子含于 \(\mathcal{FA}_{s}\) 。

(ζ) 对 A-模的每个正合列 \(0 \to \mathbf{M} \xrightarrow{h} \mathbf{N} \xrightarrow{p} \mathbf{P} \to 0\) ，序列 \(0 \longrightarrow \mathbf{M}_{\mathcal{F}} \xrightarrow{h_{\mathcal{F}}} \mathbf{N}_{\mathcal{F}} \xrightarrow{p_{\mathcal{F}}} \mathbf{P}_{\mathcal{F}} \longrightarrow 0\) 是正合的。

（证明 \((\alpha) \Rightarrow (\beta) \Rightarrow (\gamma) \Rightarrow (\alpha)\) 。为看出 \((\gamma)\) 蕴含 \((\alpha)\) ，注意对每个 A-模 M 及每个 \(x \in \mathcal{FM}\) ，存在 \(s \in S\) 使 \(sx = 0\) 。为看出 \((\beta)\) 蕴含 \((\gamma)\) ，先证明 \((\beta)\) 蕴含 \((\delta)\) 、\((\varepsilon)\) 与 \((\zeta)\) 。为证明 \((\beta)\) 蕴含 \((\zeta)\) ，利用 \((\delta)\) 确立：若 \(m \in \mathcal{F}\) 且

\[
u \in \operatorname{Hom} _ {\mathbf {A}} (\mathfrak {m}, \mathrm{P} / \mathscr {F P}),
\]

则存在 \(s \in S\) 使 \(As \subset m\) ，并且存在 \(v \in \operatorname{Hom}_{\mathbf{A}}(As, N/\mathcal{F}N)\) 使图表

\[
\begin{array}{c} \text {As} \xrightarrow {v} \text {N / FN} \\ i \Big \downarrow \qquad \qquad \qquad \qquad \qquad \qquad \qquad \qquad \Big \downarrow p ^ {\prime} \\ \mathfrak {m} \xrightarrow [ u ]{} \text {P / FP} \end{array}
\]

交换，其中 i 是典范单射，\(p'\) 是由 p 取商得到的映射。最后，对所有 \(s \in S\) ，考虑正合列

\[
0 \longrightarrow \mathrm{N} \longrightarrow \mathrm{A} \xrightarrow {\mu_ {s}} \mathrm{A} \longrightarrow \mathrm{A} / \mathrm{As} \longrightarrow 0
\]

其中 \(\mu_s\) 是位似映射 \(\lambda \mapsto \lambda s\) ；利用 \((\varepsilon)\) 、\((\zeta)\) 以及《代数》第 I 章 § 8 习题 4。）

\begin{enumerate}
\def\labelenumi{(\alph{enumi})}
\setcounter{enumi}{2}
\item
  由 (b) 推出，若 A 有一个分母在 S 中的左分式环 \((\mathbf{B}, \phi)\) ，则这个环具有下列泛性质：对每个环同态 \(f: A \to C\) （它使 \(f(s)\) 对所有 \(s \in S\) 在 C 中可逆），存在唯一的同态 \(g: B \to C\) 使 \(f = g \circ \phi\) 。特别地，B 典范地同构于 \(A_{F}\) 。
\item
  A 的分母在 S 中的右分式环的概念类似地定义。由 (c) 推出，若 A 的左分式环与右分式环都存在，则这两个环同构。
\item
  设 E 是域 K 上的一个向量空间，\((e_{t})_{t\in I}\) 是 E 的一个基，T 是 E 的张量代数。设 J 是 I 的一个异于 I 的非空子集，S 是由那些 \(e_{t}\) （满足 \(t\in J\) ）生成的 T 的乘性子集。证明 T 没有分母在 S 中的左分式环或右分式环。
\item
  沿用 (e) 的记号，假设 I = N ，考虑 T 对由元素 \(e_{0}^{2}\) 、\(e_{0}e_{i} - e_{i+1}e_{0}\) （对所有 i \textgreater{} 0）以及 \(e_{i}e_{j} - e_{j}e_{i}\) （对 i \textgreater{} 0, j \textgreater{} 0）所生成的双边理想取商所得的商环 A。设 \(S'\) 是由诸 \(e_{i}\) （对 i \textgreater{} 0）在 A 中的典范像生成的乘性集；证明 A 有分母在 \(S'\) 中的左分式环，但没有右分式环。
\item
  假设 A 有一个分母在 S 中的左分式环。在集合 \(S \times A\) 中，令 R 为 \((s, x)\) 与 \((t, y)\) 之间的关系：「存在 \(r \in S\) , \(u \in A\) , \(v \in A\) 使 r = us = vt 且 ux = vy」。证明 R 是一个等价关系，并在集合 \(B = (S \times A)/R\) 上定义一个环结构，使得由 S 与映射 \(\phi: A \to B\) （它把 \(x \in A\) 映为 \((1, x)\) 的类）所组成的序偶是 A 的一个分母在 S 中的左分式环（利用这样的事实：对 \(s \in S\) 与 \(t \in S\) ，存在 \(u \in A\) 与 \(v \in A\) 使 r = us = tv 属于 S）。考虑 \(S = A - \{0\}\) 的情形（参见《代数》第 I 章 §9 习题 8）。
\end{enumerate}

\begin{enumerate}
\def\labelenumi{\arabic{enumi}.}
\setcounter{enumi}{22}
\tightlist
\item
  \begin{enumerate}
  \def\labelenumii{(\alph{enumii})}
  \tightlist
  \item
    设 A 是一个没有零因子的左 Noether 环。证明 A 有一个左分式域（利用习题 14 (a) 证明习题 22 (b) 的条件 \((\beta_{1})\) 对 \(S = A - \{0\}\) 成立）。
  \end{enumerate}
\end{enumerate}

\begin{enumerate}
\def\labelenumi{(\alph{enumi})}
\setcounter{enumi}{1}
\tightlist
\item
  设 A 是一个没有零因子的左且右 Noether 环。证明每个有限生成的左 A-模 M（其每个 \(\neq 0\) 的元素都是自由的）都是一个自由 A-模的子模。（把 M 典范地嵌入 \(\mathbf{K} \otimes_{\mathbf{A}} \mathbf{M}\) ，其中 K 是看作右 A-模的 A 的（左或右）分式域；若 \((x_{i})\) 是左向量 K-空间 \(\mathbf{K} \otimes_{\mathbf{A}} \mathbf{M}\) 的一个基，证明 A 中存在 \(a \neq 0\) 使 M 含于 \(\mathbf{K} \otimes_{\mathbf{A}} \mathbf{M}\) 的由诸 \(a^{-1}x_{i}\) 生成的子 A-模。）
\end{enumerate}

¶ 24. 设 A 是一个环，交换与否均可，F 是 A 的满足 \(l \cap m \neq 0\) （对 A 的每个左理想 \(m \neq 0\) ）的左理想 l 组成的集合。

为使 \(l \in F\) ，充分必要条件是：对 A 的每个元素 \(a \neq 0\) ，存在 \(b \in A\) 使 \(ba \neq 0\) 且 \(ba \in l\) 。F 的每个理想 l 都包含 A 的左基座（《代数》第 VIII 章 §5 习题 9），而且若 A 是 Artin 环，则 F 由包含 A 的左基座的那些理想组成。

\begin{enumerate}
\def\labelenumi{(\alph{enumi})}
\tightlist
\item
  证明 \(\mathcal{F}\) 是一个拓扑化集（习题 16）。由此推出，集合 \(\mathfrak{a}\) ——由 \(a \in \mathbf{A}\) 中满足 \(\mathfrak{l}a = 0\) （对某个 \(\mathfrak{l} \in \mathcal{F}\) ）的那些 a 组成（诸 \(\mathfrak{l} \in \mathcal{F}\) 的右零化子的并）——是 \(\mathbf{A}\) 的一个不含幂等元的双边理想。若 \(\mathbf{A}\) 是左 Noether 环，证明 \(\mathfrak{a}\) 是幂零的（先证明每个元素 \(a \in \mathfrak{a}\) 都是幂零的，办法是考虑诸 \(a^n\) 的左零化子；然后利用《代数》第 VIII 章 § 6 习题 26 (b)）。
\end{enumerate}

设 K 是一个交换域，B 是系数在 K 中、未定元为 Y 与 \(X_{i}\) ( \(i \geqslant 1\) ) 的多项式环。设 b 是 B 的由诸 \(X_{i}X_{j}\) ( \(i \neq j\) ) 与诸 \(Y^{i-1}X_{i}\) 生成的理想，令 A 为商环 B/b；证明在 A 中，上面定义的理想 a 包含 Y 的类，而 Y 的类不是幂零的。

\begin{enumerate}
\def\labelenumi{(\alph{enumi})}
\setcounter{enumi}{1}
\item
  证明对 A 的每个左理想 m，存在 A 的一个左理想 n 使 \(m \cap n = 0\) 且 \(m + n \in F\) 。（取 n 为在满足 \(m \cap l = 0\) 的诸左理想 l 中为极大的一个理想。）
\item
  若 (a) 中定义的双边理想 a 化为 0，则称环 A 为（左）neat 环。若如此，则集合 F 是幂等的（设 l、m、n 是 A 的左理想，满足 \(n \in F\) ，\(l \subset n\) ，\(n/l\) 是 F-可忽略的，且 \(m \neq 0\) 。若 \(x \neq 0\) 属于 \(m \cap n\) ，则存在 \(r \in F\) 使 \(r.x \subset l \cap m\) ）。
\end{enumerate}

¶ 25. 沿用习题 24 的记号，假设 A 是一个左 neat 环。

\begin{enumerate}
\def\labelenumi{(\alph{enumi})}
\item
  证明典范映射 \(j_{\mathbf{A}}: \mathbf{A} \to \mathbf{A}_{\mathcal{F}}\) 是单射的，这使我们可以把 A 等同于 \(\mathbf{A}_{\mathcal{F}}\) 的一个子群。设 \(\mathcal{F}'\) 是左理想 \(l'\) （ \(\mathbf{A}_{\mathcal{F}}\) 的，满足 \(l'_{\mathcal{F}} = \mathbf{A}_{\mathcal{F}}\) ）组成的集合；证明 \(\mathcal{F}'\) 是左理想 \(l'\) （ \(\mathbf{A}_{\mathcal{F}}\) 中满足 \(l' \cap m' \neq 0\) （对每个左理想 \(m' \neq 0\) ， \(\mathbf{A}_{\mathcal{F}}\) 的））组成的集合。（先证明对每个左理想 \(m' \neq 0\) （ \(\mathbf{A}_{\mathcal{F}}\) 的），有 \(A \cap m' \neq 0\) ；另一方面注意诸理想 \(l' \in \mathcal{F}'\) 恰是 \(\mathbf{A}_{\mathcal{F}}\) 的包含某个理想 \(m \in \mathcal{F}\) 的左理想。）证明环 \(\mathbf{A}_{\mathcal{F}}\) 是左 neat 环。
\item
  证明 \(\mathbf{A}_{\mathcal{F}}\) 是一个内射 A-模，也就是说（《代数》第 II 章 §2 习题 11），对 A 的每个左理想 I 以及所有
\end{enumerate}

\[
u \in \operatorname{Hom} _ {\mathbf {A}} (\mathfrak {l}, (\mathrm{A} _ {\mathcal {F}}) _ {s}),
\]

存在 \(a' \in \mathbf{A}_{\mathcal{F}}\) 使 \(u(x) = xa'\) （对所有 \(x \in \mathfrak{l}\) ）。（利用习题 24 (b) 与习题 19 (c)。）由此推出，A-模 \(\mathbf{A}_{\mathcal{F}}\) 的每个自同态都形如 \(u: x \mapsto xa'\) ，其中 \(a' \in \mathbf{A}_{\mathcal{F}}\) （注意由习题 19 (e)，\(\operatorname{Hom}_{\mathbf{A}}(\mathbf{A}, \mathbf{A}_{\mathcal{F}}) = \operatorname{Hom}_{\mathbf{A}_{\mathcal{F}}}(\mathbf{A}_{\mathcal{F}}, \mathbf{A}_{\mathcal{F}})\) ）；同样有 \(\operatorname{Ker}(u) = (\operatorname{Ker}(u))_{\mathcal{F}}\) 。

\begin{enumerate}
\def\labelenumi{(\alph{enumi})}
\setcounter{enumi}{2}
\item
  证明对每个左理想 \(\mathfrak{l}\) （ \(\mathbf{A}\) 的），\(\mathfrak{l}_{\mathcal{F}}\) 是 \(\mathbf{A}_{\mathcal{F}}\) 的一个直因子。（设 \(\mathfrak{m}\) 是 \(\mathbf{A}\) 的一个左理想，满足 \(\mathfrak{m} \cap \mathfrak{l} = 0\) 且 \(\mathfrak{m} + \mathfrak{l} \in \mathcal{F}\) ；把 \(\mathfrak{m} + \mathfrak{l}\) 到 \(\mathfrak{m}\) 上的、在 \(\mathfrak{l}\) 上为零的投射延拓为一个自同态 \(p\) （ \(\mathbf{A}_{\mathcal{F}}\) 的）；证明 \(p^2 = p\) 且 \(\operatorname{Ker}(p) = \mathfrak{l}_{\mathcal{F}}\) 。）
\item
  证明 \(A_{F}\) 是一个绝对平坦环（第 I 章 §2 习题 17）。（化为 \(A = A_{F}\) 的情形。若 u 是自同态 \(x \mapsto xa\) （左 A-模 \(A_{s}\) 的），则 \(\operatorname{Ker}(u) = (\operatorname{Ker}(u))_{\mathcal{F}}\) ，因而 \(\operatorname{Ker}(u)\) 在 \(A_{s}\) 中有一个补理想 m；注意 \(\operatorname{Im}(u)\) 同构于 m，从而是内射的（《代数》第 II 章 §2 习题 11），因此是 \(A_{s}\) 的一个直因子（同上）。最后，利用《代数》第 II 章 §1 习题 13。）
\end{enumerate}

¶ 26. (a) 证明每个没有 ≠0 的幂零理想的 Zorn 环 A（《代数》第 VIII 章 § 6 习题 13）都是左 neat 环；特别地，每个绝对平坦环以及每个基座不化为 0 的本原环都是左 neat 环。

\begin{enumerate}
\def\labelenumi{(\alph{enumi})}
\setcounter{enumi}{1}
\item
  设 A 是一个基座不化为 0 的左本原环；证明若把 A 表示为一个向量空间的自同态环 \(\operatorname{End}_{\mathbb{D}}(\mathbf{V})\) 的稠密子环（《代数》第 VIII 章 § 5 习题 10），则与 A 相应的环 \(A_{F}\) （习题 25）同构于 \(\operatorname{End}_{\mathbb{D}}(\mathbf{V})\) 。
\item
  证明每个拟单环 A（《代数》第 VIII 章 § 5 习题 5）都是左且右 neat 环，并且相应的环 \(\mathbf{A}_{\mathcal{F}}\) 是拟单的。由此推出，对每个不化为 0 的环 B，存在一个到环 R 的同态 \(\phi: \mathrm{B} \to \mathrm{R}\) ，其中 R 是拟单的、绝对平坦的，并且 \(\mathbb{R}_s\) 是内射 R-模。
\item
  设 A 是一个没有幂零理想的左 Noether 环。证明 A 是左 neat 环，并且相应的环 \(A_{F}\) 是一个半单环（「Goldie 定理」；利用习题 24 (a)，另一方面利用《代数》第 VIII 章 § 2 习题 7，注意在绝对平坦环 \(A_{F}\) 中不可能存在任何两两正交的幂等元族，除非 A 含有左理想 \(\neq0\) 的一个无穷直和）。证明若 \(s\in A\) 在 A 中不是右零因子，则它在 \(A_{F}\) 中可逆，反之亦然；这些元素组成的集合 S 是 A 的一个乘性子集。证明 F 是 A 的与 S 相交的左理想组成的集合（为看出 A 的与 S 相交的理想属于 F，注意 S 的元素在 A 中不是左零因子；为看出每个 \(l\in F\) 都包含 S 的一个元素，把 \(A_{F}\) 的每个单分量表示为一个 n 维向量空间 E 的自同态环，并用归纳法定义 n 个元素 \(u_{1},\ldots,u_{n}\) ，使 \(u_{i}u_{j}=0\) （对 j \textless i）且 \(u_{i}^{2}\neq0\) ，诸 \(u_{i}\) 的秩为 1；最后证明 \(s=u_{1}+\cdots+u_{n}\) 的核为零）。由此推出，\((\mathbf{A}_{\mathcal{F}},j_{\mathrm{A}})\) 是 A 的一个分母在 S 中的左分式环（习题 22）。
\end{enumerate}

\begin{enumerate}
\def\labelenumi{\arabic{enumi}.}
\setcounter{enumi}{26}
\tightlist
\item
  \begin{enumerate}
  \def\labelenumii{(\alph{enumii})}
  \tightlist
  \item
    设 A 是一个环，M 是一个有限生成的忠实 A-模。证明若 a、b 是 A 的两个理想且 aM = bM，则 a 与 b 的根相等（利用第 2 节命题 4 的推论 2）。
  \end{enumerate}
\end{enumerate}

\begin{enumerate}
\def\labelenumi{(\alph{enumi})}
\setcounter{enumi}{1}
\item
  由 (a) 推出，若在整环 A 中 p 是一个有限生成的素理想，则 \(p^{m} \neq p^{n}\) （对 \(m \neq n\) ）。
\item
  设 K 是一个域，A 是两个未定元上的多项式环 K{[}X, Y{]} 。设 a 是理想 \(AX^{4} + AX^{3}Y + AXY^{3} + AY^{4}\) ，b 是 A 的理想 \(a + AX^{2}Y^{2}\) ；证明 \(b \neq a\) 且 \(b^{2} = a^{2}\) 。
\end{enumerate}
% semantic-fragments: legacy-input-seam
\relax{}
